% Chapter 6 -- Point-Cloud Completion: Method and Debugging Story
% Standalone, compilable draft (FPT-template article class + booktabs + float).
% Folded into the thesis main document later via \input.
%
% Scope: this chapter is the completion METHOD and the debugging narrative that
% produced it. The completion RESULTS (donor metric, box metric, held-out
% replication, movers) live in Chapter 7; here we only state each fix's headline
% effect and defer the full evaluation.
%
% Honesty-phrasing rules applied (personal_agenda P1):
%   - domain-gap failures (#15-19) and Step 1c (#45) presented as FINDINGS, not
%     apologies; #26 inference bug presented as a debugging lesson.
%   - final path = per-car length estimate ON, Step 1c flags OFF.
%   - 23% fallback (#41) stated as a live limitation.
%   - measured vs inferred tiers marked; numbers traceable to findings/artifacts.
%
% Figures (all embedded via \graphicspath from the gitignored output/ tree;
% regenerable from the named scripts):
%   Fig 6.1 fig61_partial_compare.py    -> fig61_partials/       (naive vs KITTI-like partial)
%   Fig 6.2 fig62_box_overlay.py        -> fig62_deblob/box_overlay_deblob_08.png
%           (BEV old vs corrected completion box vs amodal GT; reconstructs the
%            retired PCA+partial-radius normalization, touches nothing frozen)
%   Fig 6.3 viz_completion.py           -> completion_shape_08.png (final output, output/08)
% Floats set [htbp]. Draft-grade renders; may get a cosmetic restyle in the figure pass.
%
\section{Object completion}
\label{ch:completion-method}

This part of the method takes each detected car and \emph{completes} it:
from the one-sided, partial cluster a single LiDAR frame observes, it predicts
the full car surface. This section describes the completion method. It is written
as the debugging story that actually produced it, because the useful scientific
content here is not the network --- a standard Point Completion Network (PCN)
taken off the shelf --- but the sequence of failures and fixes needed to make a
completion model trained on synthetic CAD data work on real LiDAR. Several of
those steps are negative results, and they are reported as findings rather than
smoothed over. The evaluation of the finished method --- whether it adds
genuinely unseen surface, and whether that generalises --- is the subject of
Chapter~\ref{ch:completion-eval}; here we state each fix's headline effect and defer the rest.

\subsection{Completion targets}
\label{sec:completion-what}

Completion runs on the tracked car detections of Section~\ref{sec:meth-detection}. The input is a
\emph{single-frame} cluster in the sensor frame, not the accumulated multi-frame
track. This is a deliberate choice with an evidence basis (the derisk study
behind Finding~\#26): accumulating a moving car over many frames produces a
motion smear up to several metres long, which is neither a valid completion
target nor matchable to a car shape. A single frame is one clean, if partial,
observation. The value proposition --- that completing this partial improves the
car's amodal box and recovers surface the frame never saw --- is what Chapter~\ref{ch:completion-eval}
measures; this section is how the completion is produced.

\subsection{The completion network}
\label{sec:completion-pcn}

We use PCN (Yuan et al., 2018), an encoder--decoder for point clouds
(\texttt{src/pcn.py}; Figure~\ref{fig:pcn}). The encoder is a stacked PointNet: shared per-point MLPs
lift each point to a local feature, a max-pool summarises the whole cloud into a
global vector, that global vector is concatenated back onto every point, and a
second MLP-plus-max-pool stage produces a single \num{1024}-dimensional shape
embedding. The decoder is coarse-to-fine: a fully-connected stage predicts
\num{1024} coarse seed points, then a \emph{grid-folding} stage attaches a small
2D grid patch to each seed and learns to fold it into local surface geometry,
densifying the coarse skeleton into the final cloud. Training minimises a
two-term Chamfer loss, $\mathrm{CD}_{\text{coarse}} + 0.5\,\mathrm{CD}_{\text{fine}}$;
the fine term is down-weighted so the coarse skeleton is learned first. At
inference the partial is subsampled to \num{256} points before it enters the
encoder.

\begin{figure}[H]
  \centering
  \begin{tikzpicture}[
      font=\footnotesize, node distance=6mm and 7mm,
      blk/.style={draw, rounded corners, align=center, fill=orange!12, thick,
        minimum height=8mm, inner sep=3pt},
      io/.style={draw, ellipse, align=center, fill=black!5, inner sep=2pt},
      arr/.style={-{Stealth[length=2mm]}, thick},
      lbl/.style={font=\footnotesize\itshape, text=black!60},
    ]
    \node[io] (in) {partial\\\num{256} pts};
    \node[blk, right=of in] (ppm) {per-point\\MLP};
    \node[blk, right=of ppm] (mp1) {max-pool,\\concat back};
    \node[blk, right=of mp1] (mp2) {MLP,\\max-pool};
    \node[blk, right=of mp2] (emb) {shape code\\\num{1024}-d};
    \node[blk, below=10mm of emb] (fc) {FC};
    \node[blk, left=of fc] (coarse) {\num{1024}\\coarse seeds};
    \node[blk, left=of coarse] (fold) {grid-folding\\{\scriptsize 2D patch / seed}};
    \node[io, left=of fold] (out) {dense car\\\num{4096} pts};
    \draw[arr] (in)--(ppm); \draw[arr] (ppm)--(mp1); \draw[arr] (mp1)--(mp2);
    \draw[arr] (mp2)--(emb); \draw[arr] (emb)--(fc); \draw[arr] (fc)--(coarse);
    \draw[arr] (coarse)--(fold); \draw[arr] (fold)--(out);
    \node[lbl, above=1mm of mp1] {encoder (stacked PointNet)};
    \node[lbl, below=1mm of coarse] {decoder (coarse-to-fine)};
  \end{tikzpicture}
  \caption{The completion network (PCN, \texttt{src/pcn.py}). A stacked-PointNet
    encoder maps the \num{256}-point partial to a \num{1024}-d global shape code;
    a coarse-to-fine decoder predicts \num{1024} seed points and folds a small 2D
    grid patch onto each to produce the dense \num{4096}-point car. Every block is
    learned; the architecture is used off the shelf and unmodified.}
  \label{fig:pcn}
\end{figure}

The architecture is not a contribution and was not modified: a controlled
benchmark against PoinTr (Finding~\#28) found PoinTr's synthetic edge small
(Chamfer \SI{0.161}{}$\to$\SI{0.153}{\meter}) and its real-data behaviour
equivalent, so PCN was kept and architecture changes were ruled out. Everything
that follows is about the \emph{data} PCN is trained on and the \emph{geometry}
of how a real partial is presented to it --- which is where all the difficulty
turned out to live.

\subsection{Domain-gap failures (a negative result)}
\label{sec:completion-domaingap}

The first several attempts to run PCN on real cars all produced ``blobs'' ---
rounded, car-sized point clouds with no recognisable car structure (Findings
\#15--\#19). At the time this was attributed to the training data being
synthetic, and the response was a series of sim-to-real remedies: LiDAR-noise
augmentation of the synthetic partials, and real-data fine-tuning. None of them
fixed the blobs, and the fine-tuning attempts (\#16/\#17/\#19) are recorded as a
negative result: on this data, fine-tuning PCN on mined real pairs did not
recover real-car geometry. The reason it could not is the subject of the next two
subsections --- the problem was in two places at once, and neither was the network
weights.

\subsection{Synthesising KITTI-like partials}
\label{sec:completion-kittilike}

The first real cause was that the synthetic \emph{partials} PCN trained on were
structurally out-of-distribution from the real clusters it met at inference. A
derisk study comparing candidate partial generators against real seq-08 clusters
(\texttt{scratchpad/derisk\_synth\_vs\_real.py}) isolated three mismatches: the
old generator sampled a dense pinhole depth render, skipped the
\SI{0.05}{\meter} voxelisation and ground removal the real input has been
through, and added range-proportional noise (up to \SI{0.15}{\meter} at
\SI{30}{\meter}) that washed out the scan-ring banding real Velodyne data shows.

The fix (\texttt{\_render\_kitti\_like} in \texttt{src/train\_pcn.py},
Finding~\#26) makes the synthetic partial match the real one: it orients the
ShapeNet car to gravity, casts a single HDL-64E viewpoint (64 elevation beams,
\SI{0.09}{\degree} horizontal resolution) from a sensor at \SI{1.73}{\meter} and
a sampled range of \SIrange{8}{30}{\meter}, adds small \emph{absolute} noise
(\SI{0.015}{\meter}), voxelises at \SI{0.05}{\meter} to match the pipeline, and
removes points within \SI{0.30}{\meter} of the ground plane. The dense complete
mesh sample remains the target. PCN retrained on these partials
(\texttt{pcn\_kitti\_best.pth}, best val loss \num{0.1246}) is in-distribution
with the real input, and on a held-out synthetic set produces clean cars, not
blobs (Chamfer \SI{0.16}{\meter}, F-score@\SI{0.1}{\meter} \num{0.76}).
Figure~\ref{fig:kitti-like-partial} contrasts the naive and KITTI-like partials
for the same ShapeNet car.

\begin{figure}[htbp]
  \centering
  \includegraphics[width=\linewidth]{fig61_partials/partials_2668f2d9b8a2daad95551e0d9e50cb0d.png}
  \caption{The training-input distribution fix (Finding~\#26): one ShapeNet car in
    three views, with its complete-surface target (left, grey) and two candidate
    training \emph{inputs}. The naive pinhole partial (centre, orange) is dense and
    uniform; the KITTI-like partial (right, blue) reproduces what real LiDAR
    delivers --- scan-ring banding, single-viewpoint coverage, \SI{0.05}{\meter}
    voxelisation, and the ground cut. Neither is a completion output; ``better''
    here means closer to deployment, not more complete. Training on the KITTI-like
    distribution is what removes the blobs on real data.}
  \label{fig:kitti-like-partial}
\end{figure}

\subsection{Correcting the inference-normalisation error}
\label{sec:completion-inferencebug}

Retraining alone did not de-blob \emph{real} clusters, and this is the more
important lesson. The residual failure was not in the model but in the
\emph{inference code} that normalised a real partial before handing it to the
model (Finding~\#26). PCN is trained on shapes expressed in a fixed canonical car
frame, centred on the true car centre and scaled by the true car radius. The
original \texttt{complete()} instead aligned each real partial with a 3D PCA
frame and normalised it by the \emph{partial's} own centroid and radius --- a
normalisation the model had never seen. This mismatch alone produced blobs, and
it did so for \emph{every} checkpoint, including in-distribution ones (it worsened
Chamfer by roughly \num{3.5}$\times$ even on synthetic data where the model is
otherwise clean). The train/inference semantics had silently diverged; this
incident is the motivating example for the project's rule that any external
integration gets a written semantics map first.

The corrected inference path (\texttt{estimate\_canonical\_frame}) reconstructs
the training normalisation from a one-sided partial:

\begin{itemize}
  \item \textbf{Reorientation.} The sensor frame is rotated into the canonical
    car frame (X${}={}$width, Y${}={}$up, Z${}={}$length): gravity maps to the
    up axis, and the horizontal heading axis maps to length. Crucially there is
    \emph{no} PCA --- the 3D PCA alignment was part of the bug.
  \item \textbf{Full-car centre estimate.} A one-sided partial's bounding-box
    centre is not the car centre, so the centre is corrected on each axis: shift
    \emph{down} to undo the upward bias from the ground cut (the dominant
    correction, \SI{0.25}{\meter}); push \emph{outward} along width toward the
    occluded far side to a half-car-width prior (\SI{1.9}{\meter}); and push
    along length toward the occluded far end (Section~\ref{sec:completion-length}).
  \item \textbf{Scale.} Normalise by the maximum point-to-centre distance divided
    by a fixed scale correction (\num{1.137}), the calibrated ratio between a
    one-sided partial radius and the true car radius.
\end{itemize}

With this path the same real seq-08 clusters recover the correct pose and scale
(Figure~\ref{fig:deblob}): the box fitted to the old completion is mis-oriented
and mis-scaled off the amodal ground-truth box, while the box fitted to the
corrected completion aligns with it, lifting BEV box IoU at every input density
(for example \num{0.47}$\to$\num{0.85} on the dense car shown). Viewed in the
canonical car frame, the corrected completions are recognisably car-shaped: they
recover the car footprint in bird's-eye view and a rough side silhouette
(roofline included) even from very sparse input --- a clean car footprint from a
\num{64}-point cluster in the extreme case. We are nonetheless careful about the
strength of the claim: crispness varies with input density (cleaner on
well-observed cars, blobbier on the sparsest), and a completion is a filled
\num{4096}-point cloud rather than a CAD-clean surface. The sharpest completions
are the \emph{in-distribution synthetic} ones of
Section~\ref{sec:completion-kittilike} (Chamfer \SI{0.16}{\meter}); on real data
the fix's value is corrected geometry --- pose, scale, and the surface-coverage
gains quantified in Chapter~\ref{ch:completion-eval}. A calibration study on
synthetic cars (Finding~\#26) further identified the residual dominant error as
\emph{centre estimation}, not scale --- which is what motivates the geometry work
in the rest of this section.

\begin{figure}[htbp]
  \centering
  \includegraphics[width=\linewidth]{fig62_deblob/box_overlay_deblob_08.png}
  \caption{The inference-normalisation fix (Finding~\#26) in bird's-eye view at
    three input densities; \emph{same checkpoint}, only the normalisation differs.
    Each panel fits an oriented box to the completed cloud and compares it to the
    amodal ground-truth box (black) over the raw points (grey). The old PCA +
    partial-radius box (red) is mis-oriented and mis-scaled --- the 3D PCA picks an
    arbitrary principal axis rather than gravity --- while the corrected
    canonical-frame box (green) aligns with the GT, lifting BEV box IoU at every
    density: \num{0.32}$\to$\num{0.73} (sparse), \num{0.25}$\to$\num{0.74} (mid),
    \num{0.47}$\to$\num{0.85} (dense) --- per-panel values from the reconstruction
    render (\texttt{fig62\_box\_overlay.py}), not an aggregate finding. The figure
    measures \emph{pose and scale}
    via the box; the completed cloud is car-shaped in the canonical frame but
    filled rather than crisp (Chapter~\ref{ch:completion-eval} quantifies real-data surface coverage).}
  \label{fig:deblob}
\end{figure}

\subsection{L-shape input gating}
\label{sec:completion-lshape}

Even with correct normalisation, some clusters cannot complete into a plausible
car because the \emph{input} is not one car: HDBSCAN fragments (part of a car) and
merges (two cars, or a \SI{90}{\degree} heading error) both feed the completion
network implausible input. To detect them we fit the cluster's ground-plane footprint with a
search-based L-shape fit (Zhang et al., 2017), which scores edge adherence over
candidate yaw angles and stays well-defined on near-square footprints where PCA
degenerates. The fitted footprint gives an input-quality gate (Finding~\#27): a
cluster is rejected as a fragment if its fitted length is below \SI{2.7}{\meter},
or as a merge if its fitted width exceeds \SI{2.3}{\meter}.

The gate was validated on seq-08 (\texttt{scratchpad/ab\_heading.py}): of \num{47}
completed car tracks, \emph{every} implausible completion originated from a
fragment or merge input, and gating them raised completion precision from
\num{18}/\num{47} (\SI{38}{\percent}) to \num{18}/\num{26} (\SI{69}{\percent}).
An honest secondary finding from the same study: the L-shape heading itself is
neutral against plain PCA on real data --- the fit earns its place by
\emph{detecting bad inputs}, not by giving a better heading.

\subsection{Per-car length estimation}
\label{sec:completion-length}

The canonical-frame centre estimate re-centres the occluded lateral side and the
ground cut, but a partial truncated at the occluded far \emph{end} is normalised
around its observed near portion, so PCN's output stops short of the far end. The
donor metric of Chapter~\ref{ch:completion-eval} confirmed the far end as the single worst-covered
region. Two successive fixes address it, both inference-only (no retraining).

\paragraph{The fixed length prior (Finding~\#35).} A longitudinal, extend-only
push mirrors the width prior along the length axis: the centre is pushed toward
the occluded far end so the completion reaches an assumed full car length of
\SI{4.14}{\meter} (the real amodal-GT median). This more than doubled far-end
coverage on the donor metric and reversed the box-length regression (signed
$\Delta L$ \SI{-0.44}{}$\to$\SI{-0.32}{\meter}). A larger \SI{4.5}{\meter} prior
was rejected because it broke the hallucination guard by over-extending genuinely
compact cars.

\paragraph{The per-car estimate (Finding~\#36).} A single \SI{4.14}{\meter}
prior is a population median, so it is wrong in both tails --- it over-extends
compact cars and under-extends long ones. It is replaced by a per-car length,
\texttt{track\_length\_estimate}: the \nth{90} percentile of the car's own
gate-passed footprint-length fits, plus a \SI{0.12}{\meter} bias correction,
falling back to the fixed prior when a track has fewer than \num{5} gate-passed
frames. The estimator was chosen offline against amodal GT with no PCN inference:
a length-from-width shortcut is impossible (the correlation between true length
and width is \num{0.018}), and the aggregate must be a \emph{quantile}, not the
maximum, because a single contaminated L-shape fit sets the maximum. This is a
per-track quantity, so the pipeline (which owns the track) computes it and passes
it in; the completion module itself stays stateless. It removes the compact
over-extension (compact signed $\Delta L$ \SI{+0.295}{}$\to$\SI{+0.063}{\meter})
and improves the box metric overall (BEV IoU \num{0.747}$\to$\num{0.771}), at a
measured cost in pooled donor coverage that Chapter~\ref{ch:completion-eval} reports in full.

\paragraph{Limitation (Finding~\#41).} The fallback is not rare: on the
regenerated seq-08 output, \num{119} of \num{518} completed tracks
(\SI{23}{\percent}) fell back to the fixed prior because they had fewer than five
gate-passed frames. This is above a comfortable threshold and is stated as a live
limitation, not a solved problem; it is revisited in the discussion.

\subsection{The second under-extension mechanism}
\label{sec:completion-step1c}

Targeting each car's \emph{true} length still leaves completions short on normal
and long cars, which pointed to a second under-extension mechanism inside the
completion itself: PCN under-fills its normalised frame, and because the centre
push also drives the normalisation radius, the length prior doubles as an output
rescale. This was investigated under pre-registration (Finding~\#45, T13) with
two decoupled knobs: separating the radius from the length push, and a
length-axis fill factor calibrated on synthetic true GT. The change fixed the
target --- normal-car box $|\Delta L|$ \num{0.329}$\to$\num{0.232} --- but
over-extended compacts and failed the pre-registered compact non-regression guard
on both test sequences. The verdict was \textbf{not to adopt it}: the compensation is
irreducibly length-dependent, which is out of scope. The knobs remain in the code
defaulted off, and the production completion is unchanged. This negative result is
reported because it bounds the method: the residual far-end shortfall on long cars
is a known, characterised limit, not an oversight.

\subsection{The final completion path}
\label{sec:completion-summary}

To summarise, a single-frame car cluster is completed by: fitting its footprint
and rejecting fragments and merges (L-shape gate); reorienting it into the
canonical car frame with no PCA; estimating the full-car centre (down-shift for
the ground cut, width prior for the occluded side, per-car length estimate for
the occluded far end); normalising by the scale-corrected radius; running the
KITTI-like-trained PCN; and rotating the output back to the sensor frame. Every
step is inference-time geometry except the one-off retraining on KITTI-like
partials; no step fine-tunes on the evaluation cars.

Figure~\ref{fig:completion-shape} shows what this path produces on real seq-08
cars across the input-density range, in the canonical car frame. The completed
cloud (green) fills the car footprint and a rough side silhouette that the
single-frame partial (blue) leaves one-sided --- recognisably car-shaped down to a
\num{64}-point input, with fitted dimensions in the expected car range
($L \approx \SI{4.2}{\meter}$, $W \approx \SI{1.9}{\meter}$,
$H \approx \SI{1.6}{\meter}$ on the well-observed cars). Consistent with the
caveats above, the sparsest inputs complete more coarsely and can inflate the
width (Chapter~\ref{ch:completion-eval}, Finding~\#44), and the clouds are filled rather than CAD-clean;
the point here is qualitative shape, with the quantitative real-data assessment ---
whether the completion adds genuinely unseen surface, and whether that survives a
pre-registered held-out replication --- being what Chapter~\ref{ch:completion-eval} sets out to measure.

\begin{figure}[htbp]
  \centering
  \includegraphics[width=\linewidth]{completion_shape_08.png}
  \caption{What the final completion path produces on real seq-08 cars, four
    tracks spanning the input-density range (\num{4679} down to \num{64} input
    points), each in the canonical car frame: bird's-eye view (left of each pair)
    and side elevation (right). Blue is the single-frame partial input, green the
    PCN completion. The completion recovers the full car footprint and a rough
    roofline silhouette from a one-sided partial; crispness and dimensional
    accuracy degrade on the sparsest inputs (the \num{64}-point car shows the
    width inflation discussed in Chapter~\ref{ch:completion-eval}). Rendered by
    \texttt{scratchpad/viz\_completion.py} from the final-config \texttt{output/08}.}
  \label{fig:completion-shape}
\end{figure}

\section{A reference-free completion metric}
\label{sec:donor-metric}

Evaluating the completion output requires a metric, and on real LiDAR there is no
complete ground-truth surface to score against. This section explains why the
obvious choice --- Chamfer distance to a pseudo ground truth --- is invalid, then
defines the leakage-free \emph{donor-frame} metric used in its place. The metric is
validated and applied to the completion output in Chapter~\ref{ch:completion-eval};
it is defined on static cars by construction, with moving cars treated separately in
Section~\ref{sec:results-movers}.

\subsection{The invalidity of Chamfer distance on real cars}
\label{sec:chamfer-invalid}

The natural way to score a completed car is Chamfer distance (CD) against a
ground-truth surface. Real LiDAR offers no such surface: a car is seen once, from
one side. The obvious substitute is a \emph{pseudo}-ground-truth cloud --- the
labelled points of one static car accumulated across every ego frame that
observed it. This subsection explains why that substitute is not a valid target,
and hence why the completion work of Section~\ref{ch:completion-method} needed a different metric
altogether.

The fundamental limitation is that the accumulation is itself one-sided. A parked car is
observed from a road running past one of its flanks, so every frame sees broadly
the same near faces --- the road-facing side, the near end, the roofline edge ---
and none sees the occluded far side or the underbody. Pooling those frames
densifies the \emph{visible} surface but adds almost no genuinely new surface: the
pseudo-ground-truth is a denser partial, not a whole car.

Scored against such a reference, plain Chamfer \emph{rewards under-completion}.
The distance is symmetric, so every predicted point pays the distance to its
nearest reference point. A completion that correctly reconstructs the occluded
far side places points where the reference has nothing, and each such point adds
cost; a model that predicts only the already-observed surface pays nothing for
what it omits. The raw partial --- which by construction sits exactly on observed
surface and predicts nothing new --- therefore scores the \emph{lowest} (best) CD
of any cloud, on every real example (Finding~\#26). A metric that ranks the
un-completed input first cannot be used to measure completion.

This is a measured failure rather than a conjecture: across the real seq-08 cars
we scored, the raw partial had the lowest pseudo-ground-truth CD of the raw /
mirrored / completed clouds in every case (renders in
\texttt{output/experiments/verify\_pcn\_step2/}). Its consequence is that when the geometry work
began, completion had no valid real-data metric at all --- the in-distribution
synthetic CD of Section~\ref{ch:completion-method} (\SI{0.16}{\meter}) measures the network on synthetic
cars but says nothing about a real cluster. The next subsection turns the very
one-sidedness that breaks Chamfer into the basis of a metric that works.

\subsection{The donor-frame occluded-side principle}
\label{sec:donor-principle}

The idea turns the one-sidedness of LiDAR from a problem into the measurement.
For a static car observed across many frames, we complete it from \emph{one}
frame's pipeline cluster and then score the completed cloud only against surface
that the \emph{other} frames --- the \emph{donor} frames --- saw but the input
frame did not (Figure~\ref{fig:donor-principle}). Two properties make this a fair
test of completion:

\begin{itemize}
  \item The raw partial covers \emph{none} of that donor-only surface by
    construction, so any coverage the completion achieves is genuinely added,
    previously-unseen surface --- not a rescoring of the input points.
  \item The reference is built from ground-truth labelled points only, never
    from a pipeline output, so no detection or completion error can leak into
    the thing we measure against. The metric is a \emph{paired} comparison of
    raw versus completed on identical inputs, so it tolerates reference noise
    the same way the box evaluation does.
\end{itemize}

\begin{figure}[htbp]
  \centering
  \includegraphics[width=0.86\linewidth]{donor_principle.png}
  \caption{The donor-frame occluded-side principle (schematic; a car seen from
    above, the road below). The input frame observes the near flank (blue); as the
    ego drives past, the \emph{donor} frames additionally observe the two ends
    (amber), which together form the \emph{novel set} --- the surface the metric
    scores against. The far flank (grey) is occluded from every frame and is not
    part of the reference. Because the novel set lies at least $\tau$ from every
    input point, the raw partial covers none of it by construction. Illustrative
    geometry, not measured data; real per-car coverage on seq~08 is
    Figure~\ref{fig:donor-coverage}.}
  \label{fig:donor-principle}
\end{figure}

\paragraph{Relation to existing reference-free evaluation.} Completion on real LiDAR has
been scored before without a complete target, but not against the previously-unseen
surface specifically. The KITTI completion metrics introduced with
PCN~\cite{yuan2019pcnpointcompletionnetwork} --- Fidelity, Minimal Matching Distance, and
Consistency --- reference the input, a synthetic car prior, and the model's own cross-frame
outputs respectively, so a completion that merely reproduces its input can score well on
all three. Ren et al.~\cite{ren2022selfsupervised} come closest to the construction used
here: they too accumulate each tracked vehicle across frames to form a pseudo ground truth.
The difference is what the accumulation is scored against. They compare the completed cloud
to the accumulation \emph{in full}, mixing the near surface the input already observed with
the small amount the other views add; the donor metric keeps only the novel set --- the
surface at least $\tau$ from every input point --- so already-observed surface contributes
nothing and the raw partial scores zero by construction. This is also what separates the
metric from semantic scene completion~\cite{Behley_2019_ICCV}, which scores voxel occupancy
over a whole scene rather than surface coverage of one object. The novel-set restriction,
the mirrored baseline of Section~\ref{sec:donor-definition}, and the per-band hallucination
guard of Section~\ref{sec:donor-guard-lesson} are what make the metric measure
occluded-surface recovery specifically, rather than overall resemblance to an accumulated
cloud.

\subsection{Definition}
\label{sec:donor-definition}

For each (frame, instance) true-positive pair on the well-observed static cars,
we form three point sets and score them against a common target.

\paragraph{Inputs.} The \textbf{input} is the production pipeline's cluster for
that car in that frame, completed through the production path (the KITTI-like
PCN checkpoint behind the L-shape input gate). The \textbf{donor cloud} is the
union of the instance's ground-truth (labelled-car) points from every observing
frame \emph{except} the input frame, transformed to the world frame and
voxelised at \SI{3}{\centi\meter}.

\paragraph{Novel set (the visibility mask).} From the donor cloud we keep only
points that lie at least $\tau$ from \emph{every} input point (nearest-neighbour
test). This \emph{novel set} is exactly the surface the input frame did not
observe. We use $\tau = \SI{0.15}{\meter}$ as the primary radius and
\SI{0.10}{} / \SI{0.20}{\meter} as stability checks.

\paragraph{Scored clouds.} We score three clouds against the novel set so that
completion is judged against meaningful baselines rather than in isolation:
\begin{itemize}
  \item \textbf{raw} --- the partial input (the floor: covers no novel surface
    by construction);
  \item \textbf{mirrored} --- the input unioned with its own reflection across
    the car's estimated symmetry plane (a cheap, hallucination-free way to add
    surface; the baseline completion must beat to justify a learned model);
  \item \textbf{completed} --- the PCN output.
\end{itemize}

\paragraph{Scores.} Each cloud gets three numbers, all one-directional from the
novel set to the cloud:
\begin{itemize}
  \item \textbf{coverage} (\texttt{cov@0.1}) --- the fraction of novel points
    within \SI{0.1}{\meter} of the cloud (the primary quality score: how much
    unseen surface was recovered);
  \item \textbf{median novel distance} --- the median distance from novel points
    to the cloud (a threshold-free companion to coverage);
  \item \textbf{out-of-box} --- the fraction of the cloud's points falling
    outside the amodal ground-truth box (plus a \SI{0.2}{\meter} margin): a
    \emph{hallucination guard}, since a model can trivially raise coverage by
    inventing surface everywhere.
\end{itemize}

\paragraph{Eligibility.} The metric applies to any completion whose input is a
\emph{single, well-formed object cluster} --- one cluster corresponding to one
car, with a footprint consistent with a whole vehicle --- so that there is a
coherent object shape to score. In this pipeline that admission is realised by the
L-shape input gate of Section~\ref{ch:completion-method} (fitted footprint length
$\geq \SI{2.7}{\meter}$ and width $\leq \SI{2.3}{\meter}$, rejecting fragments and
merges), but the metric itself does not depend on that particular gate: any
admission rule that isolates a single object serves in its place.

\paragraph{Statistics.} A pair counts if its input is admitted and its novel set
has at least \num{100} points at the primary $\tau$. Because the frames of one
parked car are autocorrelated, we reduce each car to its per-car median and test
across cars with the Wilcoxon signed-rank test --- the same paired, per-car design
used for the box evaluation.

